\section{Related literature}\label{sec:lit}

\subsection{Why arbitrage is limited}
The idea that arbitrage keeps prices in line rests on the arbitrageur being able to trade without limit and without
risk. Shleifer and Vishny (1997) point out that arbitrage is done by specialists who invest other people's money. When a
mispricing widens, the arbitrageur shows losses and investors withdraw funds, so capital leaves exactly when the
opportunity is best. De Long, Shleifer, Summers and Waldmann (1990) show that the risk created by noise traders can by
itself keep prices away from fundamentals, because an arbitrageur with a short horizon may have to close at a worse
price. Gromb and Vayanos (2010) bring these mechanisms into one framework. Their account has three features that the
model of this report shares: a demand shock moves the price by more when arbitrageurs are more risk averse or have less
capital; a cost of short selling matters only when arbitrageurs need to be short, so the band around fair value is
asymmetric; and margin requirements make capital a binding constraint. Short-sale constraints have a separate literature.
Miller (1977) argues that when opinions differ and pessimists cannot sell short, prices reflect the optimists and are
too high, and Diamond and Verrecchia (1987) show that such constraints slow the response of prices to bad news. These
predictions are the subject of the second stage of the project.

\subsection{The no-arbitrage band}
A future and the stock bought with borrowed money give the same payoff at expiry, so the futures price should equal the
cost of carry, and trading costs create a band around it within which arbitrage does not pay. Inside the band the basis
should wander; outside, it should be pushed back. This calls for a threshold model. Dwyer, Locke and Yu (1996) estimate a
three-regime threshold model for S\&P 500 index futures with minute data and find thresholds of about half a round-trip
cost and faster reversion outside the band than inside. Tse (2001) argues that arbitrageurs differ in their costs, so
the market's response should strengthen smoothly with the size of the gap, and fits a smooth-transition model to Dow
Jones futures. Ters and Urban (2020) show that in a threshold error-correction model the arbitrage cost is estimated
precisely even when its components are not, and suggest comparing the band before and after regulatory changes. I use
threshold and smooth-transition models from this tradition, applied to a daily panel of stocks rather than to a single
index.

\subsection{Evidence from India}
Indian studies of the stock futures basis are few and use short intraday samples (\cref{tab:litbands}). Aggarwal (2015)
estimates threshold models for 141 stock futures with one-second order-book data over six months of 2013 and finds an
average band of about $-15$ to $+3$\bp, wider for less liquid stocks, and moving with funding conditions. Shankar,
Sankar and Kiran Kumar (2019) estimate weekly threshold models for 102 stock futures with five-minute data from 2007 to
2009 and find that both thresholds are negative, at about $-39$ and $-10$\bp, with the lower edge more sensitive to
volatility, a sign of a constraint on the short side. Slivka, Shah and Wu (2011) add up the costs of arbitrage between a
stock and its future and find a band about twice as wide as in mature markets. At the index level, Nandan, Agrawal and
Bhargava (2014) find that Nifty futures traded below their cost-of-carry value on most days from 2003 to 2012. Saurav,
Srivastava and Agarwalla (2026) study the stock futures basis around earnings announcements and short-sale constraints.

\begin{table}[t]
\caption{Earlier Indian estimates describe intraday bands; this report measures the overnight band}\label{tab:litbands}
\exsub{Studies of single-stock futures against their underlying stocks on the National Stock Exchange}
\small\begin{tabularx}{\linewidth}{@{}K{3.4cm}P{2.9cm}P{2.4cm}L@{}}
\toprule
\thd{Study} & \thd{Sample} & \thd{Data} & \thd{Band} \\
\midrule
Aggarwal (2015) & 141 stocks, Mar to Aug 2013 & One-second order book & About $-15$ to $+3$\bp, per stock and day \\ \rowrule
Shankar, Sankar and Kiran Kumar (2019) & 102 stocks, 2007 to 2009 & Five-minute prices & $-39$ to $-10$\bp, weekly \\ \rowrule
Slivka, Shah and Wu (2011) & Cost calculation & Fee schedules & About twice the width in mature markets \\ \rowrule
This report & 360 stocks, 2015 to 2026 & Daily volume-weighted prices & $-34.5$ to $+23.9$\bp\ around each stock's median \\
\bottomrule
\end{tabularx}
\exsource{Source: the studies cited; \cref{sec:band} for this report.}
\end{table}

The intraday bands describe traders who open and close within minutes. The daily basis studied here measures the part of
the gap that lasts overnight, where an arbitrageur must carry the position and pay the taxes on delivery, so it should
be wider. To my knowledge no study has examined the F\&O ban or modelled a position limit that, once reached, lets every
participant only reduce positions. That is the gap this project fills, and this report lays the ground for it by
measuring the band and the ban's effect on trading.
